\subsection{The compact cubic block: a common curvature orientation}
\label{sharp:upper-cubic-section}

\begin{proposition}[The complete compact cubic carrier]
\label{sharp:upper-cubic}
Fix finite \(U_0,G_0\). On the coefficient block
\begin{equation}\label{sharp:upper-cubic-block}
 a=-1,\quad p=tU,\quad\gamma=t^2G,\quad s=t^3,
             \qquad0\le U\le U_0,\quad |G|\le G_0,
\end{equation}
the same displacement has at most two critical points and at most three
isolated zeros, counted with multiplicity, on its whole physical carrier
under the common positive unscaled raw cap. All captured target cycles
lie in this carrier after the common coefficient and marked
Hausdorff-collar shrink.
\end{proposition}

\begin{proof}
Put \(\tau=t^2\). Proposition~\ref{deg:prop-compact-cubic} supplies the
actual clipped carrier and its allocation \eqref{deg:compact-physical-cover}.
Its bounded-action and independent root-collar estimates give the same
curvature orientation. The growing range has positive curvature by
\eqref{deg:compact-curvature}, and is identified with the bounded
range by the actual resection
\begin{equation}\label{sharp:upper-cubic-overlap}
 p_\tau(u)=u+O_{C^3}(\tau(1+|\log\tau|)),\quad p_\tau'>0,
                \qquad2M+\beta\le u\le3M-\beta.
\end{equation}
These are estimates for the completed physical map, so
\begin{equation}\label{sharp:upper-cubic-common-crossing}
 \Pi''(x)=0\quad\Longrightarrow\quad\Pi'''(x)<0
\end{equation}
on the connected lower receiver. Real and fold fibres retain their
one-sided through endpoints.

The original ghost tail has the derivative bound
\eqref{deg:compact-raw-tail-derivative}; the core has
\eqref{deg:compact-core-variation}. The construction places
their comparison and high intervals in strict overlap:
\begin{equation}\label{sharp:upper-cubic-height-overlap}
 \eta_{\rm comp}<\tau K_H(1+\tau U^2)<\eta_{\max}/2,
                      \qquad\tau^2\eta_{\max}<h_c/2.
\end{equation}
One final multiplier bound makes \(F'>1/2\) on this complement.
The cover in Proposition~\ref{deg:prop-compact-cubic} includes every
input, retaining the original clips and equality dividers. Its whole-carrier
count, equivalently Lemma~\ref{sharp:upper-curvature-count}, gives
the stated multiplicity bounds. Capture is applied after these analytic
choices, as in Lemma~\ref{lem:loc-capture}.
\end{proof}
